\section{Commemorations and the Count of Orations}
\label{sec:commemorations}

\subsection{The two kinds}

\rn{106} makes the chapter apply to both books at once: what is here laid down about
commemorations holds for the Mass as for the Office, whether in occurrence or in concurrence.
\rn{107} divides them: commemorations are either privileged or ordinary. \rn{108} gives the
practical difference in liturgical placement — privileged commemorations are made at Lauds and
at Vespers and in all Masses; ordinary commemorations are made only at Lauds, in conventual
Masses, and in all low Masses.

The consequence for the Mass is exact and often misstated: an ordinary commemoration is not made
in a sung Mass that is not conventual. \rn{111}\,a says the same thing from the other side, and
Missal rubrics \rn{434}\,a repeat it for the orations.

\rn{109} lists the privileged commemorations exhaustively:

\begin{studybox}{\rn{109}: privileged commemorations}
(a) of a Sunday; (b) of a liturgical day of the I class; (c) of the days within the octave of
the Nativity of the Lord; (d) of the Ember ferias of September; (e) of the ferias of Advent, of
Lent, and of the Passion; (f) of the greater Litanies, in the Mass.\\[0.35em]
\latin{Omnes aliae commemorationes sunt commemorationes ordinariae} — all other commemorations
are ordinary.
\end{studybox}

Two entries repay attention. Clause (d) singles out the September Ember ferias among the three
sets of Ember days; the Advent and Lenten Ember ferias are covered instead by clause (e), which
privileges the ferias of those seasons generally. And clause (f) is confined to the Mass, in
keeping with \rn{81}'s rule that nothing is done about the greater Litanies in the Office.

\rn{110} adds the one inseparable commemoration in the system: in the Office and Mass of St Peter
a commemoration of St Paul is always made, and conversely; the two orations are held to coalesce
into one so completely that, in counting orations, they are reckoned as a single one. In the
Office the second Apostle's oration is added at Lauds and Vespers under a single conclusion to
the oration of the day, without antiphon and versicle; in the Mass, under a single conclusion to
the oration of the day; and whenever one Apostle's oration is to be added by way of
commemoration, the other is added to it immediately, before all other commemorations.

\subsection{How many commemorations a day admits}

\rn{111} is the counting rule, and it is the single provision by which the 1960 reform most
changed the daily texture of the Roman Rite.

\begingroup\small
\renewcommand{\arraystretch}{1.15}
\begin{longtable}{@{}>{\raggedright\arraybackslash}p{0.44\linewidth} >{\raggedright\arraybackslash}p{0.48\linewidth}@{}}
\toprule
\textbf{Liturgical day} & \textbf{Commemorations admitted} \\
\midrule
\endfirsthead
\toprule
\textbf{Liturgical day} & \textbf{Commemorations admitted} \\
\midrule
\endhead
\bottomrule
\endfoot
I class days, and sung Masses that are not conventual & One privileged commemoration only, and no other \\
Sundays of the II class & One only, namely of a II class feast — and even that is omitted if a privileged commemoration must be made \\
Other II class days & One only, either one privileged or one ordinary \\
III and IV class days & Two only \\
\end{longtable}
\endgroup


\rn{112} then adds four exclusions that operate regardless of the count:

\begin{enumerate}[label=(\alph*), leftmargin=1.7em, itemsep=0.25em]
\item an Office, Mass, or commemoration of a feast or mystery of one divine Person excludes the
commemoration or oration of another feast or mystery of the same divine Person;
\item an Office, Mass, or commemoration of a Sunday excludes the commemoration or oration of a
feast or mystery of the Lord, and conversely;
\item an Office, Mass, or commemoration of the season excludes another commemoration of the
season;
\item likewise, an Office, Mass, or commemoration of the Blessed Virgin Mary or of any Saint or
Blessed excludes another commemoration or oration in which the intercession of the same Blessed
Virgin, Saint, or Blessed is implored — \latin{quod tamen non valet de oratione dominicae vel
feriae, in qua fit invocatio eiusdem Sancti}, which however does not hold of the oration of a
Sunday or feria in which an invocation of the same Saint occurs.
\end{enumerate}

\rn{113} fixes the order: the commemoration of the season is made first; in admitting and
arranging the others, the order of the table of precedence is kept. \rn{114} disposes of the
remainder without ceremony: \latin{Quaelibet commemoratio, quae numerum pro singulis diebus
liturgicis statutum superet, omittitur} — any commemoration exceeding the number laid down for
the individual liturgical days is omitted.

The scale of the change is best seen against what preceded it, and the older ceiling can be
quoted. The pre-1955 witness described in Section~\ref{sec:witness} requires, of the orations
that may be added at a Mass, that they \latin{septenarium numerum non excedant, atque imparem
praeterea numerum} keep — that they not exceed the number seven, and that they keep besides an
odd number. Seven, not three, was the outer bound, and the count was required to be odd; a low
Mass could accordingly carry the Mass's own oration, the commemorations, the seasonal oration,
the \latin{Suffragium}, the commemoration of the Cross, the imperated oration, and the priest's
votive oration, and stacking three, five, or seven orations was ordinary. The 1955 decree cut
deeply; the 1960 code caps the count at the figures above and, in the Missal at \rn{435},
replaces the odd-number requirement with a flat ceiling of three.

\subsection{The Missal's implementation: orations}

Missal rubrics \rnn{433--465} restate the counting rules in terms of orations, and add the
provisions that only the Mass needs.

\rn{433} enumerates what counts as an ``oration'' in the Mass: (a) the oration of the Mass being
celebrated; (b) the orations of a commemorated Office and of any occurring commemoration;
(c) other orations prescribed by the rubrics (\rnn{447--453}); (d) the oration imperated by the
local Ordinary (\rnn{454--460}); (e) the votive oration that may be said on certain liturgical
days at the celebrant's discretion (\rnn{461--465}).

\rn{434} gives the counts, which match \rn{111} exactly and extend them to votive Masses of the
corresponding classes: on I class liturgical days, in I class votive Masses, and in sung
non-conventual Masses, no oration beyond the Mass's own is admitted except one to be said under
a single conclusion and one privileged commemoration; on II class Sundays, only the
commemoration of a II class feast, omitted if a privileged commemoration must be made; on other
II class days and in II class votive Masses, one other only, privileged or ordinary; on III and
IV class days and in III and IV class votive Masses, two only.

\rn{435} adds the absolute ceiling and states it with unusual force:

\begin{studyframe}
\latin{Quaelibet oratio, quae numerum pro singulis diebus liturgicis statutum superet, omittitur;
profecto numerum ternarium orationum nullo praetextu excedere licet.}
\medskip

\textit{Working gloss.} Any oration exceeding the number laid down for the individual liturgical
days is omitted; indeed it is on no pretext permitted to exceed the number three.
\end{studyframe}

\rnn{436--437} govern conclusions: the Mass's own oration is always said under its own
conclusion unless another oration is to be joined to it under the same conclusion; and
commemorations, the imperated oration, and the votive oration are always said under a separate
conclusion. \rn{438} handles near-duplication: if two orations are composed in nearly the same
words in their first or second part, the later one is changed — if of the season, to that of the
following Sunday or feria; if of a Saint, to another of the same or a similar Common; if it is
an imperated oration, it is omitted. \rn{439} forbids altering the words \latin{hanc} or
\latin{hodiernam} or \latin{praesentem diem} and the like in the orations of a transferred or
reposited Office — a small rule with a large implication, since it means the 1962 system
tolerates a transferred feast praying about ``this day''.

\rnn{444--446} restrict the joining of orations under a single conclusion to three cases: a
ritual oration (\rn{447}); the oration of an impeded I or II class votive Mass; and any other
oration that the rubrics expressly indicate or concede as to be said under a single conclusion
with the Mass's oration, the cross-references naming \rnn{110}, 350, 449, 451, and 453. Only one
such oration may be joined; if several would qualify, one only is kept, in the order just given,
and the rest are omitted. An oration said under a single conclusion is counted together with the
Mass's oration, and is said also in sung Masses.

\rnn{449--453} supply the three standing occasions for such an added oration: the day of the
Supreme Pontiff's coronation and its anniversary, and the anniversary of the diocesan bishop's
election, consecration, or translation, once a year on a day of the bishop's choosing
(\rn{449}); the anniversary of a priest's own ordination (\rn{451}); and the second-to-last
Sunday of October or another day appointed by the local Ordinary for the Missions, on which the
oration \latin{pro Fidei propagatione} is added (\rn{453}). All three are barred on the days
listed under numbers 1, 2, 3, and 8 of the table of precedence — that is, on Christmas, Easter,
and Pentecost; on the Sacred Triduum; on the Epiphany, Ascension, Trinity, Corpus Christi, the
Sacred Heart, and Christ the King; and on the Commemoration of All the Faithful Departed.

\rnn{454--460} govern the imperated oration, which the local Ordinary may impose when a grave
and public necessity or calamity occurs. It may be only one; it must be said by all priests
celebrating in the churches and oratories of the diocese, even exempt ones; it is never said
under a single conclusion with the Mass's oration, but after the privileged commemorations; and
it is forbidden on all I and II class liturgical days, in I and II class votive Masses, in sung
Masses, and whenever privileged commemorations have already filled the number laid down.
\rn{456} counsels that the Ordinary impose it not in a settled way but only for genuinely grave
cause and for no longer than the true necessity. \rn{459} provides for a calamity of long
duration — war, pestilence, and the like — in which case the oration is said only on Mondays,
Wednesdays, and Fridays. \rn{460} allows a parish priest, when there is no time to approach the
Ordinary, to appoint a suitable oration for three continuous days within his parish, subject to
the same bars.

\rnn{461--465} govern the votive oration, which is the celebrant's own liberty and is
correspondingly narrow: any priest may add one oration at will in all low non-conventual Masses
on IV class liturgical days. It is placed last, after the other orations, and must not carry the
total beyond three.

\subsection{What the commemoration rules suppressed}

The \latin{Variationes in Ordinario divini Officii} state one suppression in a single line, and
it is the one most often noticed by those who moved between the pre-1955 and the 1962 books:
\latin{Suffragium de omnibus Sanctis et commemoratio de Cruce abolentur} — the Suffrage of All
Saints and the commemoration of the Cross are abolished. Both were standing additions rather than
occasional ones; their removal, together with the caps of \rn{111} and the ceiling of \rn{435},
is what produced the characteristic single-oration Mass of the 1962 books.

Nothing in these provisions touches the calendar's inventory. A feast reduced to a commemoration
by the \latin{Variationes} is still inscribed in the calendar and still generates an oration when
the day's count admits one; a feast struck from the calendar generates nothing. The distinction
between the two operations is worth keeping sharp, because a reference that conflates them will
misreport both the shape of the calendar and the content of a given day's Mass.
